\section{Problem setting}\label{sec:setting}

Consider a 2-armed stochastic bandit over $T$ rounds. The reward $X_{i,t}$ from arm $i$ at round $t$ is independent and $1$-sub-Gaussian, i.e.
\[
\mathbb E\bigl[e^{\lambda(X_{i,t}-\mu_i(t))}\bigr]\le e^{\lambda^2/2},\qquad \forall\lambda\in\mathbb R,
\]
where $\mu_i(t)$ is the mean of arm $i$ at round $t$. There are two known change points $0<\tau_1<\tau_2<T$, defining three stationary phases:
\[
\mu_i(t)=\begin{cases}\mu_i, & t\le\tau_1,\\ \mu_i', & \tau_1<t\le\tau_2,\\ \mu_i'', & t>\tau_2.\end{cases}
\]
Only arm 1's mean changes at $\tau_1$ (i.e.\ $\mu_2=\mu_2'$), while both arms' means may change at $\tau_2$. Let $\Delta=|\mu_1-\mu_2|$ and $\Delta'=|\mu_1'-\mu_2'|$. Rewards are bounded in $[0,1]$; the width of this range (i.e.\ $1$) is the sensitivity referenced in part~(d).

\section{(a) A Hoeffding-type bound}\label{sec:hoeffding}

\begin{proposition}
For $n$ independent $1$-sub-Gaussian samples with common mean $\mu$, $\mathbb P(|\widehat\mu_n-\mu|\ge t)\le 2\exp(-nt^2/2)$.
\end{proposition}

\begin{proof}
Let $S_n=\sum_{j=1}^n(X_j-\mu)$ and $\widehat\mu_n=n^{-1}\sum_{j=1}^nX_j$. Independence and the sub-Gaussian assumption give, for every $\lambda\in\mathbb R$,
\[
\mathbb E[e^{\lambda S_n}]=\prod_{j=1}^n\mathbb E[e^{\lambda(X_j-\mu)}]\le e^{n\lambda^2/2}.
\]
For a deviation threshold $u>0$ and $\lambda>0$, Markov's inequality yields
\[
\mathbb P(\widehat\mu_n-\mu\ge u)=\mathbb P(e^{\lambda S_n}\ge e^{\lambda nu})\le e^{-\lambda nu}\mathbb E[e^{\lambda S_n}]\le\exp\Bigl(-\lambda nu+\frac{n\lambda^2}{2}\Bigr).
\]
Since $-\lambda nu+n\lambda^2/2=\tfrac n2(\lambda-u)^2-\tfrac{nu^2}2$, choosing $\lambda=u$ gives $\mathbb P(\widehat\mu_n-\mu\ge u)\le e^{-nu^2/2}$. Applying the same argument to $-S_n$ gives $\mathbb P(\widehat\mu_n-\mu\le-u)\le e^{-nu^2/2}$. A union bound therefore proves
\begin{equation}\label{eq:hoeffding}
\mathbb P(|\widehat\mu_n-\mu|\ge u)\le 2e^{-nu^2/2},\qquad u\ge0,
\end{equation}
where the case $u=0$ is immediate.
\end{proof}

\section{(b) UCB in the first phase}\label{sec:ucb}

Let $I_t\in\{1,2\}$ be the selected arm and $N_i(t)=\sum_{s=1}^t\mathbf 1\{I_s=i\}$. The cumulative pseudo-regret is
\[
R(a{:}b)=\sum_{t=a}^b\Bigl(\max_i\mu_i(t)-\mu_{I_t}(t)\Bigr).
\]
Write $H=\tau_1$. For $\Delta>0$, let $i^\star$ and $j$ be the optimal and suboptimal arms in the first phase. Then $\mu_{i^\star}-\mu_j=\Delta$ and $R(1{:}H)=\Delta N_j(H)$. If $\Delta=0$, the pseudo-regret is zero; if $H=1$, it is at most $\Delta$. Assume henceforth that $H\ge2$ and $\Delta>0$.

\paragraph{Concentration at adaptive counts.} Let $(\mathcal F_t)$ be the observation filtration, including the policy's independent randomization, so that $I_t$ is $\mathcal F_{t-1}$-measurable. With $\xi_t=X_{I_t,t}-\mu_{I_t}(t)$, independence implies $\mathbb E[e^{\lambda\xi_t}\mid\mathcal F_{t-1}]\le e^{\lambda^2/2}$. Fix an arm $i$ and deterministic pull indices $a+1,\dots,a+m$. Set
\[
A_s=\mathbf 1\{I_s=i,\ a<N_i(s-1)+1\le a+m\},\qquad C_t=\sum_{s\le t}A_s,\qquad S_t=\sum_{s\le t}A_s\xi_s.
\]
The indicator $A_s$ is predictable, so $\mathbb E[e^{\lambda A_s\xi_s-\lambda^2A_s/2}\mid\mathcal F_{s-1}]\le1$. Hence $M_t(\lambda)=\exp(\lambda S_t-\lambda^2C_t/2)$ is a nonnegative supermartingale with $\mathbb EM_t(\lambda)\le1$. At a deterministic horizon $h$, completion of the block means $C_h=m$. For $u,\lambda>0$, Markov's inequality gives $\mathbb P(C_h=m,\ S_h\ge mu)\le e^{-\lambda mu+m\lambda^2/2}$. Taking $\lambda=u$ and applying the same argument to $-S_h$ gives the completed-block bound
\begin{equation}\label{eq:block}
\mathbb P(C_h=m,\ |S_h|\ge mu)\le 2e^{-mu^2/2}.
\end{equation}

\paragraph{Algorithm.} Fix $\delta\in(0,1)$ and set
\[
L_H=\log\Bigl(\frac{4H}{\delta}\Bigr),\qquad r_H(n)=\sqrt{\frac{2L_H}{n}}.
\]
Pull each arm once. At the start of round $t$, let $\widehat\mu_i(t)=N_i(t-1)^{-1}\sum_{s<t,\,I_s=i}X_{i,s}$. For $3\le t\le H$, choose
\begin{equation}\label{eq:ucb}
I_t\in\arg\max_{i\in\{1,2\}}\bigl\{\widehat\mu_i(t)+r_H(N_i(t-1))\bigr\},
\end{equation}
with arbitrary tie-breaking. The gap $\Delta$ is not needed by the algorithm.

\paragraph{Regret bound.} For each arm and each prefix length $n\le H$, \eqref{eq:block} bounds the probability that the prefix is observed and violates its confidence bound by $2e^{-nr_H(n)^2/2}=2e^{-L_H}=\delta/(2H)$. A union bound over both arms and all $n\le H$ gives an event $\mathcal E_H$ of probability at least $1-\delta$ on which $|\widehat\mu_{i,n}-\mu_i|\le r_H(n)$ for every observed prefix. These bounds therefore hold at the adaptive counts $N_i(t-1)$. On $\mathcal E_H$, selection of the suboptimal arm $j$ at $t>2$ implies
\[
\mu_{i^\star}\le\widehat\mu_{i^\star}(t)+r_H(N_{i^\star}(t-1))\le\widehat\mu_j(t)+r_H(N_j(t-1))\le\mu_j+2r_H(N_j(t-1)).
\]
Thus $\Delta\le2\sqrt{2L_H/N_j(t-1)}$, i.e.\ $N_j(t-1)\le 8L_H/\Delta^2$. Apply this inequality to the last post-initialization pull of $j$ and add one for that pull. If no such pull occurs, $N_j(H)=1$. In either case $N_j(H)\le1+8L_H/\Delta^2$. Together with $N_j(H)\le H$, this yields, with probability at least $1-\delta$,
\begin{equation}\label{eq:ucb_bound}
R(1{:}H)\le\min\Bigl\{H\Delta,\ \Delta+\frac{8\log(4H/\delta)}{\Delta}\Bigr\}.
\end{equation}
Taking $\delta=H^{-2}$ gives $R(1{:}H)\le\Delta+8\log(4H^3)/\Delta$. Since $\log(4H^3)\le5\log H$ for $H\ge2$ and $0<\Delta\le1$,
\[
R(1{:}\tau_1)=O\Bigl(\frac{\log\tau_1}{\Delta}\Bigr)\quad\text{with probability at least }1-\tau_1^{-2}.
\]

\section{(c) Sliding-window estimation across the change points}\label{sec:window}

After $\tau_1$, the learner uses a sliding-window estimator: at round $t>\tau_1$, it estimates each arm's mean using only the last $w$ observed pulls of that arm. The same sliding-window mechanism continues unchanged across $\tau_2$ (there is no reset at $\tau_2$).

\subsection*{(i) Error bound after $\tau_1$}

Fix an integer $1\le w\le T$. At the start of round $t$, let $\mathcal W_i(t)$ contain the times of arm $i$'s last $m_i(t)=\min\{w,N_i(t-1)\}$ observed pulls. For $m_i(t)>0$, define
\[
\widehat\mu_i^{\mathrm{SW}}(t)=\frac1{m_i(t)}\sum_{s\in\mathcal W_i(t)}X_{i,s},\qquad \overline\mu_i(t)=\frac1{m_i(t)}\sum_{s\in\mathcal W_i(t)}\mu_i(s).
\]
Let $B_i(t)=\overline\mu_i(t)-\mu_i(t)$ be the signed mixing bias and $\mathcal B_i(t)=|B_i(t)|$. Then
\begin{equation}\label{eq:decomp}
\widehat\mu_i^{\mathrm{SW}}(t)-\mu_i(t)=\frac1{m_i(t)}\sum_{s\in\mathcal W_i(t)}\bigl(X_{i,s}-\mu_i(s)\bigr)+B_i(t).
\end{equation}
After pull $n$, the window consists of the pull-index block $\max\{1,n-w+1\},\dots,n$. There are at most $T$ such blocks per arm. Applying \eqref{eq:block} at horizon $T$ and taking a union bound over these blocks gives, with
\[
L_T=\log\Bigl(\frac{4T}{\delta}\Bigr),\qquad r(m)=\sqrt{\frac{2L_T}{m}},
\]
the simultaneous bound
\begin{equation}\label{eq:sw_bound}
|\widehat\mu_i^{\mathrm{SW}}(t)-\mu_i(t)|\le r(m_i(t))+\mathcal B_i(t)
\end{equation}
with probability at least $1-\delta$, for both arms and every nonempty window. An empty window is initialized by sampling its arm.

For $t>\tau_k$, write
\[
q_{i,k}(t)=\sum_{s=\tau_k+1}^{t-1}\mathbf 1\{I_s=i\},\qquad \alpha_{i,k}(t)=\frac{\#\{s\in\mathcal W_i(t):s\le\tau_k\}}{m_i(t)}.
\]
The window contains $\min\{m_i(t),q_{i,k}(t)\}$ post-change observations, so
\begin{equation}\label{eq:alpha}
\alpha_{i,k}(t)=\frac{(m_i(t)-q_{i,k}(t))_+}{m_i(t)},\qquad (x)_+=\max\{x,0\}.
\end{equation}
In particular, for a full window, $\alpha_{i,k}(t)=(1-q_{i,k}(t)/w)_+$.

For $\tau_1<t\le\tau_2$, put $J_1=|\mu_1'-\mu_1|$. Since $\mu_2=\mu_2'$, $\overline\mu_1(t)=\alpha_{1,1}(t)\mu_1+(1-\alpha_{1,1}(t))\mu_1'$ and $\overline\mu_2(t)=\mu_2'$. Hence $B_1(t)=\alpha_{1,1}(t)(\mu_1-\mu_1')$ and $B_2(t)=0$. Substitution into \eqref{eq:sw_bound} gives
\begin{equation}\label{eq:phase2}
|\widehat\mu_1^{\mathrm{SW}}(t)-\mu_1'|\le r(m_1(t))+\alpha_{1,1}(t)J_1,\qquad |\widehat\mu_2^{\mathrm{SW}}(t)-\mu_2'|\le r(m_2(t)).
\end{equation}
For full windows, this becomes
\begin{equation}\label{eq:phase2_full}
|\widehat\mu_1^{\mathrm{SW}}(t)-\mu_1'|\le\sqrt{\frac{2L_T}{w}}+\Bigl(1-\frac{q_{1,1}(t)}{w}\Bigr)_+J_1,\qquad |\widehat\mu_2^{\mathrm{SW}}(t)-\mu_2'|\le\sqrt{\frac{2L_T}{w}}.
\end{equation}
The jump $J_1$ governs the mixing bias, whereas $\Delta'=|\mu_1'-\mu_2'|$ is the phase-2 inter-arm gap. Since $\mu_2=\mu_2'$, the triangle inequality gives $J_1=|(\mu_1'-\mu_2)-(\mu_1-\mu_2)|\le\Delta'+\Delta$. Thus a bound using only the two stated gaps replaces the bias term by $\alpha_{1,1}(t)(\Delta+\Delta')$; it cannot in general be replaced by $\alpha_{1,1}(t)\Delta'$.

\subsection*{(ii) Extension past $\tau_2$}

For $t>\tau_2$, let
\[
\begin{aligned}
n_{i,1}(t)&=\#\{s\in\mathcal W_i(t):s\le\tau_1\},\\
n_{i,2}(t)&=\#\{s\in\mathcal W_i(t):\tau_1<s\le\tau_2\},\\
n_{i,3}(t)&=\#\{s\in\mathcal W_i(t):\tau_2<s<t\}.
\end{aligned}
\]
Suppressing $t$ in the counts and writing $m_i=m_i(t)$, $n_{i,1}+n_{i,2}+n_{i,3}=m_i$ and $\overline\mu_i(t)=(n_{i,1}\mu_i+n_{i,2}\mu_i'+n_{i,3}\mu_i'')/m_i$. Subtracting the current mean yields
\begin{equation}\label{eq:bias3}
B_i(t)=\frac{n_{i,1}}{m_i}(\mu_i-\mu_i'')+\frac{n_{i,2}}{m_i}(\mu_i'-\mu_i'').
\end{equation}
Set $d_{i,1}=\mu_i'-\mu_i$, $d_{i,2}=\mu_i''-\mu_i'$ and $J_{i,2}=|d_{i,2}|$. Since $\mu_i-\mu_i''=-d_{i,1}-d_{i,2}$ and $\mu_i'-\mu_i''=-d_{i,2}$, and $\alpha_{i,1}=n_{i,1}/m_i$, $\alpha_{i,2}=(n_{i,1}+n_{i,2})/m_i$, equation \eqref{eq:bias3} becomes
\begin{equation}\label{eq:bias3b}
B_i(t)=-\alpha_{i,1}(t)d_{i,1}-\alpha_{i,2}(t)d_{i,2}.
\end{equation}
Using $d_{2,1}=0$ in \eqref{eq:sw_bound} gives
\begin{equation}\label{eq:phase3}
|\widehat\mu_1^{\mathrm{SW}}(t)-\mu_1''|\le r(m_1(t))+\alpha_{1,1}(t)J_1+\alpha_{1,2}(t)J_{1,2},\qquad
|\widehat\mu_2^{\mathrm{SW}}(t)-\mu_2''|\le r(m_2(t))+\alpha_{2,2}(t)J_{2,2}.
\end{equation}
For full windows,
\begin{equation}\label{eq:phase3_full}
\begin{aligned}
|\widehat\mu_1^{\mathrm{SW}}(t)-\mu_1''|&\le\sqrt{\frac{2L_T}{w}}+\Bigl(1-\frac{q_{1,1}(t)}{w}\Bigr)_+J_1+\Bigl(1-\frac{q_{1,2}(t)}{w}\Bigr)_+J_{1,2},\\
|\widehat\mu_2^{\mathrm{SW}}(t)-\mu_2''|&\le\sqrt{\frac{2L_T}{w}}+\Bigl(1-\frac{q_{2,2}(t)}{w}\Bigr)_+J_{2,2}.
\end{aligned}
\end{equation}
For arm 1, the first-change term in \eqref{eq:bias3b} must be retained while phase-1 observations remain in its window. For arm 2, this term is identically zero. The second-change term involves all observations from phases 1 and 2. The reward range also gives the jump-independent bound
\begin{equation}\label{eq:range_bias}
\mathcal B_i(t)\le\frac{n_{i,1}}{m_i}|\mu_i-\mu_i''|+\frac{n_{i,2}}{m_i}|\mu_i'-\mu_i''|\le\alpha_{i,2}(t).
\end{equation}
Once $q_{i,2}(t)\ge w$, the window contains only phase-3 observations and its mixing bias vanishes. Renewal is therefore measured in post-change pulls of the corresponding arm, not in calendar rounds.

\subsection*{(iii) One window for both change points?}

A common window can be suitable when the two phases have comparable precision and renewal requirements. In general, its choice depends on the post-change gaps, jump magnitudes and available fresh samples, not only on whether one or two arms change. The asymmetry nevertheless favours different memory lengths in a more flexible design: retaining arm 2's history at $\tau_1$, and refreshing both potentially changed arms at $\tau_2$. With the fixed-window mechanism in the question, these requirements must be balanced by a common $w$.

Let $g>0$ be the current inter-arm gap. On the confidence event,
\[
\bigl|(\widehat\mu_1^{\mathrm{SW}}(t)-\widehat\mu_2^{\mathrm{SW}}(t))-(\mu_1(t)-\mu_2(t))\bigr|\le r(m_1(t))+r(m_2(t))+|B_1(t)-B_2(t)|.
\]
An upper bound strictly below $g$ guarantees correct empirical ranking. For full windows, reserving half the gap for sampling fluctuation gives
\begin{equation}\label{eq:wprec}
2\sqrt{\frac{2L_T}{w}}\le\frac g2\iff w\ge\frac{32L_T}{g^2}.
\end{equation}
It then suffices that $|B_1(t)-B_2(t)|<g/2$. By contrast, for a jump $J>0$ and a tolerance $0\le\eta<J$,
\begin{equation}\label{eq:renewal}
J(1-q/w)_+\le\eta\iff q\ge w(1-\eta/J).
\end{equation}
Indeed, for $q<w$ this follows by rearranging $1-q/w\le\eta/J$; for $q\ge w$, the bias is zero. Thus larger windows improve precision but require proportionally more fresh pulls to reduce a substantial mixing bias.

For nonzero post-change gaps, the sufficient statistical scales are $L_T/(\Delta')^2$ and $L_T/(\Delta'')^2$, with $\Delta''=|\mu_1''-\mu_2''|$. When they differ greatly, choosing $w$ for the smaller gap can lead to unnecessarily slow renewal in the other phase. When they are comparable and both phases provide enough fresh pulls of the affected arms, one window may perform well at both changes. Since each suboptimal pull incurs the current gap, prolonging recovery can directly increase regret.

A two-arm shift need not be more difficult for ranking. If phase-1 observations have disappeared, $d_{1,2}=d_{2,2}=d$, and both stale fractions equal $\alpha$, then $B_1=B_2=-\alpha d$. This common bias cancels in the estimated gap. The jump directions and the sampling proportions must therefore enter the choice of $w$. The sampling policy must also ensure fresh observations of each affected arm, since its window does not change while it is unselected.

\section{(d) Laplace observation noise and a privacy budget}\label{sec:noise}

Suppose each reward at round $t$ is only observed after adding independent noise $Z_t\sim\mathrm{Laplace}(0,b_t)$, at a cost $\varepsilon_t=c/b_t$ per round for a known constant $c>0$ tied to the sensitivity above. The learner chooses a (possibly time-varying) noise schedule $\{b_t\}_{t=1}^T$ subject to a fixed total budget $\varepsilon_{\mathrm{total}}=\sum_{t=1}^T\varepsilon_t$. Write
\[
Y_t=X_{I_t,t}+Z_t,\qquad \varepsilon_t=\frac c{b_t},\qquad E=\varepsilon_{\mathrm{total}}>0,\qquad \sum_{t=1}^T\varepsilon_t=E.
\]
The scale is chosen before the observation. Conditional on the past and this scale, $Z_t$ is a fresh Laplace draw independent of the current reward. The cost uses the sensitivity of a single raw reward, whose range has width $1$.

\paragraph{Confidence bound.} First consider independent observations $Y_j=X_j+Z_j$, $1\le j\le n$, with common reward mean $\mu$ and deterministic scales $b_j$. For $Z\sim\mathrm{Laplace}(0,b)$, direct integration gives
\[
\mathbb E[e^{\lambda Z}]=\frac1{2b}\int_{\mathbb R}e^{\lambda z-|z|/b}\,dz=\frac1{1-b^2\lambda^2},\qquad|\lambda|<1/b.
\]
In particular, $\mathbb EZ=0$ and $\mathrm{Var}(Z)=2b^2$. The moment-generating function diverges for $|\lambda|\ge1/b$, so the noisy observations are not globally sub-Gaussian. Put $S_n=\sum_{j=1}^n(Y_j-\mu)$, $V_n=n+2\sum_{j=1}^nb_j^2$ and $b_*=\max_jb_j$. For $|\lambda|<1/b_*$, independence and the sub-Gaussian assumption imply
\begin{equation}\label{eq:subgamma}
\log\mathbb E[e^{\lambda S_n}]\le\frac{n\lambda^2}2-\sum_{j=1}^n\log(1-b_j^2\lambda^2)\le\frac{V_n\lambda^2}{2(1-b_*|\lambda|)}.
\end{equation}
The last inequality follows from $b_j\le b_*$ and $-\log(1-u^2)=\sum_{k\ge1}u^{2k}/k\le u^2/(1-u^2)\le u^2/(1-u)$ for $0\le u<1$. For $x>0$, take $a=\sqrt{2x/V_n}$ and $\lambda=a/(1+b_*a)$. Then
\[
-\lambda\bigl(\sqrt{2V_nx}+b_*x\bigr)+\frac{V_n\lambda^2}{2(1-b_*\lambda)}=-\frac{2x+b_*ax}{1+b_*a}+\frac{x}{1+b_*a}=-x.
\]
The exponential Markov inequality, applied to $S_n$ and $-S_n$, therefore gives
\begin{equation}\label{eq:bernstein}
\mathbb P\bigl(|S_n|\ge\sqrt{2V_nx}+b_*x\bigr)\le2e^{-x}.
\end{equation}
For the sample mean, \eqref{eq:bernstein} gives the confidence radius $(\sqrt{2V_n\ell}+b_*\ell)/n$, where $\ell=\log(2/\delta)$. For a constant scale $b_j=b$, this is
\begin{equation}\label{eq:rb}
r_b(n,\delta)=\sqrt{\frac{2(1+2b^2)\log(2/\delta)}n}+\frac{b\log(2/\delta)}n.
\end{equation}
Thus Laplace noise adds a variance contribution and a linear tail term; replacing the variance parameter alone does not give the full bound.

For the noisy sliding-window estimator $\widetilde\mu_i^{\mathrm{SW}}(t)=m_i(t)^{-1}\sum_{s\in\mathcal W_i(t)}Y_s$, the bias $B_i(t)$ is unchanged from part~(c). For constant $b$, the conditional exponential-moment argument used in \eqref{eq:block}, followed by a union bound over the windows, gives
\begin{equation}\label{eq:noisy_sw}
|\widetilde\mu_i^{\mathrm{SW}}(t)-\mu_i(t)|\le\mathcal B_i(t)+\sqrt{\frac{2(1+2b^2)L_T}{m_i(t)}}+\frac{bL_T}{m_i(t)}
\end{equation}
simultaneously with probability at least $1-\delta$, where $L_T=\log(4T/\delta)$. For a predetermined schedule taking $K$ distinct scales, define
\[
V_i(t)=m_i(t)+2\sum_{s\in\mathcal W_i(t)}b_s^2,\quad b_{i,*}(t)=\max_{s\in\mathcal W_i(t)}b_s,\quad D_T=1+\lceil\log_2T\rceil,\quad L_*=\log\Bigl(\frac{4TKD_T}\delta\Bigr).
\]
The corresponding simultaneous bound is
\begin{equation}\label{eq:noisy_sw_tv}
|\widetilde\mu_i^{\mathrm{SW}}(t)-\mu_i(t)|\le\mathcal B_i(t)+\frac{2\sqrt{V_i(t)L_*}+b_{i,*}(t)L_*}{m_i(t)}.
\end{equation}
To obtain it, apply the conditional form of \eqref{eq:subgamma} to each pull-index block, stopping before any observation with scale greater than a candidate maximum $\beta$. For a block of length $m$ with this maximum, $m+2\beta^2\le V\le m(1+2\beta^2)$; at most $D_T$ dyadic ranges suffice. Within each range, an upper endpoint $v$ satisfies $V\le v\le2V$, which replaces $\sqrt{2vx}$ by $2\sqrt{Vx}$. The union bound over at most $2T$ windows, $K$ maxima and $D_T$ ranges has failure probability at most $4TKD_Te^{-x}$; take $x=L_*$. Since $K\le T$, $L_*=O(\log(T/\delta))$.

\paragraph{Choosing $w$ and the schedule.} Uniform expenditure gives $\varepsilon_t=E/T$ and $b_t=cT/E$, so the statistical radius of a full window is
\begin{equation}\label{eq:uniform}
\sqrt{\frac{2[1+2(cT/E)^2]L_T}w}+\frac{cTL_T}{Ew}.
\end{equation}
A smaller budget increases the noise scale and the averaging required for a given precision. Increasing $w$ reduces this statistical error but raises the renewal requirement in \eqref{eq:renewal}. For a current gap $g>0$, let $r_i(t)$ be the statistical radius in \eqref{eq:noisy_sw} or \eqref{eq:noisy_sw_tv}. A sufficient empirical-ranking condition is
\begin{equation}\label{eq:ranking}
r_1(t)+r_2(t)+\mathcal B_1(t)+\mathcal B_2(t)<g.
\end{equation}
For full windows with constant scale $b$, requiring each of the two terms in the radius to be at most $g/8$ gives
\[
\sqrt{\frac{2(1+2b^2)L_T}w}\le\frac g8\iff w\ge\frac{128(1+2b^2)L_T}{g^2},\qquad \frac{bL_T}w\le\frac g8\iff w\ge\frac{8bL_T}g.
\]
Each arm then has statistical radius at most $g/4$. Thus the sufficient window scale is
\begin{equation}\label{eq:wnoise}
w\gtrsim\frac{(1+b^2)L_T}{g^2}+\frac{bL_T}g,
\end{equation}
with a sufficiently large universal constant, together with total mixing bias below $g/2$. Equations \eqref{eq:wnoise} and \eqref{eq:renewal} determine the precision--adaptation trade-off: more accurate recovery observations can support a smaller window and shorten the period of uncertain arm choices.

A concrete budget-feasible family is obtained by partitioning the horizon into initialization, post-change recovery and stable blocks $\mathcal I_\ell$ of lengths $h_\ell$. Choose positive allocations $E_\ell$ with $\sum_\ell E_\ell=E$, and set
\begin{equation}\label{eq:blocks}
\varepsilon_t=\frac{E_\ell}{h_\ell},\qquad b_t=\frac{ch_\ell}{E_\ell},\qquad t\in\mathcal I_\ell.
\end{equation}
Then $\sum_t\varepsilon_t=\sum_\ell h_\ell(E_\ell/h_\ell)=E$. A predetermined partition and allocation yield a schedule covered by \eqref{eq:noisy_sw_tv}.

The window and block allocations should be chosen jointly to satisfy the target precision in \eqref{eq:ranking} while limiting renewal costs. At fixed recovery-block length, increasing its allocation reduces the noise of fresh observations. More expenditure is useful when resolving uncertainty can improve decisions over many remaining rounds. Recovery after $\tau_1$ should prioritize fresh pulls of arm 1; recovery after $\tau_2$ must sample both potentially changed arms. The relative allocations depend on the gaps, jump magnitudes and remaining phase lengths.

Stable blocks must retain enough accuracy for the chosen fixed window. If their noise scales are too large, later observations replace earlier accurate ones and inflate both $\sum_{s\in\mathcal W_i(t)}b_s^2$ and $b_{i,*}(t)$. Budget concentration around changes is therefore balanced against maintaining reliable comparisons throughout the stable periods.
